\documentclass{beamer}

\usepackage[T2A]{fontenc}
\usepackage[utf8]{inputenc}
\usepackage[english,russian]{babel}
\input{settings.tex}


\title{Матричные неравества, Оптимальное управление, Робастность}
\subtitle{Вычислительный интеллект, Лекция 9}
\author{Сергей Савин}
\centering
\date{\mydate}



\begin{document}
\maketitle


%\begin{frame}{Content}
%	
%	\begin{itemize}
%		\item LMI
%		\item Control design
%		\item Robustness check
%		\item S-procedure
%		\item Robust Control Design
%		\item Quadratic stability
%		\item Appendix A
%	\end{itemize}
%	
%	
%\end{frame}




\begin{frame}{Синтез управления (непрерывное время), 1}
	\begin{flushleft}
		
		Рассмотрим систему $\dot{\bo{x}} = \bo{A}\bo{x} + \bo{B}\bo{u}$, управление $\bo{u} = \bo{K}\bo{x}$ и функцию Ляпунова $V = \bo{x}^\top\bo{S}\bo{x}$, $\bo{S} \succ 0$.
		
		\bigskip
		
		Замкнутая система имеет вид $\dot{\bo{x}} = (\bo{A} + \bo{B}\bo{K})\bo{x}$, а полная производная функции Ляпунова:
		
		\begin{equation}
			\dot V = \bo{x}^\top (\bo{A} + \bo{B}\bo{K})^\top\bo{S}\bo{x} + \bo{x}^\top\bo{S} (\bo{A} + \bo{B}\bo{K}) \bo{x} < 0
		\end{equation}
		
		Это можно переписать как ЛМН:
		
		\begin{equation}
			\label{eq:vdot}
			(\bo{A} + \bo{B}\bo{K})^\top\bo{S} + \bo{S} (\bo{A} + \bo{B}\bo{K}) \prec 0
		\end{equation}
		
		Это \emph{не линейно} относительно переменных решения ($\bo{S}$ и $\bo{K}$) и не может быть решено напрямую с помощью популярных решателей.
		
	\end{flushleft}
\end{frame}

\begin{frame}{Синтез управления (непрерывное время), 2}
	\begin{flushleft}
		
		Вводя новую переменную $\bo{P} = \bo{S}^{-1} \succ 0$ и умножая \eqref{eq:vdot} на $\bo{P}$ с обеих сторон (это можно сделать, так как $\bo{P}$ и $\bo{S}$ имеют полный ранг, и это конгруэнтное преобразование, сохраняющее определённость, см. приложение), получаем:
		
		\begin{equation}
			\bo{P}(\bo{A} + \bo{B}\bo{K})^\top + (\bo{A} + \bo{B}\bo{K})\bo{P} \prec 0
		\end{equation}
		
		Теперь введём ещё одну переменную $\bo{L} = \bo{K}\bo{P}$ и получим ЛМН ограничение $\bo{P}\bo{A}^\top + \bo{A}\bo{P} + \bo{L}^\top\bo{B}^\top + \bo{B}\bo{L} \prec 0$. Тогда синтез управления сводится к задаче полуопределённого программирования (SDP):
		%
		\begin{equation}
			\label{control_design}
			\begin{aligned}
				& \underset{\bo{P}, \bo{L}}{\text{minimize}} 
				& &  0, \\
				& \text{subject to:} & & 
				\begin{cases}
					\bo{P}\bo{A}^\top + \bo{A}\bo{P} + \bo{L}^\top\bo{B}^\top + \bo{B}\bo{L} \prec 0, \\ 
					\bo{P} \succ 0
				\end{cases}
			\end{aligned}
		\end{equation}
		
		Решение \eqref{control_design} даёт нам $\bo{P}$ и $\bo{L}$, из которых мы можем вычислить $\bo{K} = \bo{L}\bo{P}^{-1}$ и $\bo{S} = \bo{P}^{-1}$, решая исходную задачу.
		
	\end{flushleft}
\end{frame}

\begin{frame}{Условие Ляпунова для дискретного времени}
	\begin{flushleft}
		
		Дискретная динамическая система $\bo{x}_{i+1} = \bo{f}(\bo{x}_i)$ устойчива, если существует функция Ляпунова $V(\bo{x}_i)$, такая что:
		
		\begin{itemize}
			\item $V(\bo{x}_{i+1}) - V(\bo{x}_i) < 0$;
			\item $V(\bo{x}_i) > 0$ для всех $\bo{x}_i \neq 0$.
		\end{itemize}
		
	\end{flushleft}
\end{frame}

\begin{frame}{Синтез управления (дискретное время), 1}
	\begin{flushleft}
		
		Рассмотрим систему $\bo{x}_{i+1} = \bo{A}\bo{x}_i + \bo{B}\bo{u}_i$, управление $\bo{u}_i = \bo{K}\bo{x}_i$ и функцию Ляпунова $V(\bo{x}_i) = \bo{x}_i\T \bo{S} \bo{x}_i$, $\bo{S} \succ 0$.
		
		\bigskip
		
		Замкнутая система имеет вид $\bo{x}_{i+1} = (\bo{A} + \bo{B}\bo{K})\bo{x}_i$, а дискретная динамика функции Ляпунова:
		
		\begin{align}
			V(\bo{x}_{i+1}) - V(\bo{x}_i) <0 
			\\
			\bo{x}_{i+1}\T \bo{S} \bo{x}_{i+1} - \bo{x}_i\T \bo{S} \bo{x}_i < 0
			\\
			\bo{x}_i\T(\bo{A} + \bo{B}\bo{K})\T \bo{S} (\bo{A} + \bo{B}\bo{K})\bo{x}_i - \bo{x}_i\T \bo{S} \bo{x}_i < 0
			\\
			\bo{x}_i\T((\bo{A} + \bo{B}\bo{K})\T \bo{S} (\bo{A} + \bo{B}\bo{K}) - \bo{S})\bo{x}_i < 0
		\end{align}
		
		Последнее уравнение эквивалентно следующему матричному неравенству:
		%
		\begin{align}
			(\bo{A} + \bo{B}\bo{K})\T \bo{S} (\bo{A} + \bo{B}\bo{K}) - \bo{S} \prec 0
		\end{align}
		
		
	\end{flushleft}
\end{frame}

\begin{frame}{Синтез управления (дискретное время), 2}
	\begin{flushleft}
		
		Матричное неравенство $(\bo{A} + \bo{B}\bo{K})\T \bo{S} (\bo{A} + \bo{B}\bo{K}) - \bo{S} \prec 0$ можно переписать:
		%
		\begin{align}
			\label{discrete_LMI_1}
			-\bo{S} - (\bo{A} + \bo{B}\bo{K})\T (-\bo{S}) (\bo{A} + \bo{B}\bo{K})  \prec 0
		\end{align}
		
		Определим $\bo{P}^{-1} = \bo{S} \succ 0$. Мы можем умножить \eqref{discrete_LMI_1} на $\bo{P}$ с обеих сторон (конгруэнтное преобразование):
		%
		\begin{align}
			-\bo{P} - \bo{P}(\bo{A} + \bo{B}\bo{K})\T (-\bo{P}^{-1}) (\bo{A} + \bo{B}\bo{K})\bo{P}  \prec 0
			\\
			-\bo{P} - (\bo{A}\bo{P} + \bo{B}\bo{K}\bo{P})\T (-\bo{P}^{-1}) (\bo{A}\bo{P} + \bo{B}\bo{K}\bo{P})  \prec 0
		\end{align}
		
		Определим $\bo{L} = \bo{K}\bo{P}$:
		%
		\begin{align}
			-\bo{P} - (\bo{A}\bo{P} + \bo{B}\bo{L})\T (-\bo{P}^{-1}) (\bo{A}\bo{P} + \bo{B}\bo{L})  \prec 0
		\end{align}
		
	\end{flushleft}
\end{frame}


\begin{frame}{Дополнение Шура}
	\begin{flushleft}
		
		\begin{theorem}[Дополнение Шура]
			Пусть дана матрица $\bo{M} = \begin{bmatrix}
				\bo{A} & \bo{B} \\
				\bo{B}\T & \bo{C}
			\end{bmatrix}$, где $\bo{C} \prec 0$. 
			Следующие утверждения эквивалентны:
			\begin{enumerate}
				\item $\bo{A} - \bo{B} \bo{C}^{-1}\bo{B}\T \prec 0$
				\item $\bo{M} \prec 0$
			\end{enumerate}
		\end{theorem}
		
		
		
	\end{flushleft}
\end{frame}


\begin{frame}{Синтез управления (дискретное время), 3}
	\begin{flushleft}
		
		Матричное неравенство $-\bo{P} - (\bo{A}\bo{P} + \bo{B}\bo{L})\T (-\bo{P}^{-1}) (\bo{A}\bo{P} + \bo{B}\bo{L})  \prec 0$ можно преобразовать с помощью дополнения Шура:
		%
		\begin{align}
			\begin{bmatrix}
				-\bo{P} & \bo{A}\bo{P} + \bo{B}\bo{L} \\
				(\bo{A}\bo{P} + \bo{B}\bo{L})\T & -\bo{P}
			\end{bmatrix}  \prec 0
		\end{align}
		
		Это линейное матричное неравенство. Задача синтеза управления принимает вид:
		%
		\begin{equation}
			\label{control_design_discrete}
			\begin{aligned}
				& \underset{\bo{P}, \bo{L}}{\text{minimize}} 
				& &  0, \\
				& \text{subject to:} & & 
				\begin{cases}
					\begin{bmatrix}
						-\bo{P} & \bo{A}\bo{P} + \bo{B}\bo{L} \\
						(\bo{A}\bo{P} + \bo{B}\bo{L})\T & -\bo{P}
					\end{bmatrix}  \prec 0, \\ 
					\bo{P} \succ 0
				\end{cases}
			\end{aligned}
		\end{equation}
		
		
		
	\end{flushleft}
\end{frame}

\begin{frame}{Робастность, 1}
	\begin{flushleft}
		
		Рассмотрим систему $\dot{\bo{x}} = \bo{A}\bo{x}$, в которой матрица $\bo{A}$ известна неточно. Это не значит, что нам ничего не известно о модели, но в наших знаниях присутствует неопределённость.
		
		\bigskip
		
		Это можно смоделировать следующим образом:
		
		\begin{equation}
			\label{eq:uncertain}
			\dot{\bo{x}} = (\bo{A} + \bo{F} \Delta \bo{E})\bo{x}
		\end{equation}
		%
		где $\bo{F}$ и $\bo{E}$ — известные матрицы, а $\Delta$ — неопределённая матрица с \emph{ограниченной нормой}: $||\Delta|| \leq 1$: множество всех возможных моделей, которые мы можем ожидать, ограничено возможными значениями $\Delta$.
		
	\end{flushleft}
\end{frame}

\begin{frame}{Робастность, 2}
	\begin{flushleft}
		
		Запишем уравнение Ляпунова для системы \eqref{eq:uncertain}:
		
		\begin{equation}
			\dot V = \bo{x}^\top 
			(\bo{A} + \bo{F} \Delta \bo{E})^\top\bo{S}\bo{x} + \bo{x}^\top\bo{S} (\bo{A} + \bo{F} \Delta \bo{E}) \bo{x} < 0
		\end{equation}
		
		Введём новую переменную $\bo{w} = \Delta \bo{E}\bo{x}$:
		
		\begin{equation}
			\label{eq:Lyapunov_xw}
			\dot V = \bo{x}^\top 
			(\bo{A}^\top \bo{S} + \bo{S}\bo{A}) \bo{x} + 
			\bo{w}^\top \bo{F}^\top\bo{S} \bo{x} +
			\bo{x}^\top \bo{S}\bo{F} \bo{w} < 0
		\end{equation}
		
		Рассмотрим $\bo{w}^\top \bo{w}$:
		
		\begin{equation}
			\label{eq:Delta-inequality}
			\bo{w}^\top \bo{w} = 
			\bo{x}^\top\bo{E}^\top \Delta^\top \Delta \bo{E}\bo{x}
			\leq
			\bo{x}^\top\bo{E}^\top \bo{E}\bo{x}
		\end{equation}		
		%
		что выполняется, поскольку $|| \Delta ||\leq 1$. Фактически, единственное свойство нормы, которое нам здесь нужно, — это выполнение дельта-неравенства \eqref{eq:Delta-inequality}.
		
	\end{flushleft}
\end{frame}

\begin{frame}{Робастность, 3}
	\begin{flushleft}
		
		Имея $\bo{w}^\top \bo{w} 
		\leq
		\bo{x}^\top\bo{E}^\top \bo{E}\bo{x}$, мы можем записать:
		
		\begin{equation}
			\bo{x}^\top\bo{E}^\top \bo{E}\bo{x} - \bo{w}^\top \bo{w} \geq 0
		\end{equation}	
		
		Что эквивалентно:
		
		\begin{equation}
			\label{eq:EEww}
			\begin{bmatrix}
				\bo{x} \\ \bo{w}
			\end{bmatrix}^\top
			\begin{bmatrix}
				\bo{E}^\top \bo{E} & 0 \\
				0 & -\bo{I}
			\end{bmatrix}		
			\begin{bmatrix}
				\bo{x} \\ \bo{w}
			\end{bmatrix}
			\geq 0
		\end{equation}	
		
		Тем же способом мы можем переписать \eqref{eq:Lyapunov_xw}:
		
		\begin{equation}
			\label{eq:xw_Lyapunov}
			\begin{bmatrix}
				\bo{x} \\ \bo{w}
			\end{bmatrix}^\top
			\begin{bmatrix}
				\bo{A}^\top \bo{S} + \bo{S}\bo{A} & \bo{S}\bo{F} \\
				\bo{F}^\top\bo{S} & 0
			\end{bmatrix}
			\begin{bmatrix}
				\bo{x} \\ \bo{w}
			\end{bmatrix}
			\leq 0
		\end{equation}	
		%
		которая должна выполняться только тогда, когда выполняется \eqref{eq:EEww}.
		
	\end{flushleft}
\end{frame}

\begin{frame}{S-процедура}
	\begin{flushleft}
		
		Существует способ обеспечить выполнение ограничения $\bo{z}^\top \bo{M}\bo{z} \leq 0$ для тех $\bo{z}$, для которых $\bo{z}^\top \bo{N}\bo{z} \geq 0$. Это называется \emph{S-процедурой}.
		
		\begin{theorem}[S-процедура]
			Если $\exists \gamma \geq 0$ такое, что $\bo{M} + \gamma \bo{N} \preceq 0$, то $\bo{z}^\top \bo{N}\bo{z} \geq 0 \implies \bo{z}^\top \bo{M}\bo{z} \leq 0$ 
		\end{theorem}
		
		\bigskip
		
		Доказательство достаточности. Выберем такой $\bo{z}$, что $\bo{z}^\top \bo{N}\bo{z} \geq 0$. Тогда $0 \geq \bo{z}^\top (\bo{M} + \gamma \bo{N}) \bo{z} \geq \bo{z}^\top \bo{M}\bo{z}$. \textcolor{mygray}{(Подробнее см. в Приложении)}.
		
		
	\end{flushleft}
\end{frame}


\begin{frame}{Робастность, 4}
	\begin{flushleft}
		
		
		Используя S-процедуру, мы подтверждаем, что \eqref{eq:xw_Lyapunov} выполняется, когда выполняется \eqref{eq:EEww}:
		%
		\begin{equation}
			\begin{bmatrix}
				\bo{x} \\ \bo{w}
			\end{bmatrix}^\top
			\begin{bmatrix}
				\bo{A}^\top \bo{S} + \bo{S}\bo{A} + \gamma \bo{E}^\top \bo{E} & \bo{S}\bo{F} \\
				\bo{F}^\top\bo{S} & -\gamma\bo{I}
			\end{bmatrix}		
			\begin{bmatrix}
				\bo{x} \\ \bo{w}
			\end{bmatrix}
			\leq 0
		\end{equation}
		%
		В форме ЛМН это выглядит так:
		%
		\begin{equation}
			\begin{bmatrix}
				\bo{A}^\top \bo{S} + \bo{S}\bo{A} + \gamma \bo{E}^\top \bo{E} & \bo{S}\bo{F} \\
				\bo{F}^\top\bo{S} & -\gamma\bo{I}
			\end{bmatrix}		
			\preceq 0
		\end{equation}	
		%
		Задача полуопределённого программирования (SDP), подтверждающая робастную устойчивость для всех допустимых $\Delta$, имеет вид:
		%
		%
		\begin{equation}
			\label{robust_stability_sdp}
			\begin{aligned}
				& \underset{\bo{S}, \gamma}{\text{minimize}} 
				& &  0, \\
				& \text{subject to:} & & 
				\begin{cases}
					\begin{bmatrix}
						\bo{A}^\top \bo{S} + \bo{S}\bo{A} + \gamma \bo{E}^\top \bo{E} & \bo{S}\bo{F} \\
						\bo{F}^\top\bo{S} & -\gamma\bo{I}
					\end{bmatrix}\preceq 0, \\ 
					\bo{S} \succ 0, 
					\gamma > 0
				\end{cases}
			\end{aligned}
		\end{equation}
		
		
	\end{flushleft}
\end{frame}

\begin{frame}{Робастный синтез управления, 1}
	\begin{flushleft}
		
		Рассмотрим систему: 
		
		\begin{equation}
			\dot{\bo{x}}  = (\bo{A} + \bo{F} \Delta \bo{E})\bo{x} + \bo{B}\bo{u}
		\end{equation}
		%
		Синтезируем управление для этой системы, гарантирующее устойчивость.
		
		\bigskip
		
		Предположим, что управление линейно: $\bo{u} = \bo{K}\bo{x}$. Система принимает вид:
		%
		\begin{equation}
			\dot{\bo{x}}  = (\bo{A} + \bo{B}\bo{K} + \bo{F} \Delta \bo{E})\bo{x}
		\end{equation}
		
		Неравенство Ляпунова становится:
		%
		\begin{equation}
			\bo{x}\T(\bo{A} + \bo{B}\bo{K} + \bo{F} \Delta \bo{E})\T\bo{S}\bo{x}
			+
			\bo{x}\T \bo{S} (\bo{A} + \bo{B}\bo{K} + \bo{F} \Delta \bo{E})\bo{x} < 0
		\end{equation}
		
		
		
	\end{flushleft}
\end{frame}

\begin{frame}{Робастный синтез управления, 2}
	\begin{flushleft}
		
		Введём новую переменную $\bo{w} = \Delta \bo{E} \bo{x}$:
		
		\begin{equation*}
			\bo{x}\T(\bo{A} + \bo{B}\bo{K})\T\bo{S}\bo{x}
			+
			\bo{x}\T \bo{S} (\bo{A} + \bo{B}\bo{K})\bo{x}
			+
			\bo{x}\T\bo{S}\bo{F}\bo{w} + \bo{w}\T\bo{F}\T\bo{S}\bo{x} < 0
		\end{equation*}
		
		Перепишем это в матричной форме:
		
		\begin{equation}
			\label{eq:quad_xw_1}
			\begin{bmatrix}
				\bo{x} \\ \bo{w}
			\end{bmatrix}\T
			\begin{bmatrix}
				(\bo{A} + \bo{B}\bo{K})\T\bo{S} + \bo{S} (\bo{A} + \bo{B}\bo{K}) & \bo{S}\bo{F} \\
				\bo{F}\T\bo{S} & 0
			\end{bmatrix}
			\begin{bmatrix}
				\bo{x} \\ \bo{w}
			\end{bmatrix} < 0
		\end{equation}		
		
		Это уравнение должно выполняться только для $\bo{w} = \Delta \bo{E} \bo{x}$.
		
		
		
	\end{flushleft}
\end{frame}

\begin{frame}{Робастный синтез управления, 3}
	\begin{flushleft}
		
		
		Заметим, что $\bo{w}\T\bo{w} \leq \bo{x}\T \bo{E}\T \bo{E} \bo{x}$, поскольку $|| \Delta || \leq 1$. В матричной форме это:
		%
		\begin{equation}
			\begin{bmatrix}
				\bo{x} \\ \bo{w}
			\end{bmatrix}\T
			\begin{bmatrix}
				\bo{E}\T \bo{E} & 0 \\
				0 & -\bo{I}
			\end{bmatrix}
			\begin{bmatrix}
				\bo{x} \\ \bo{w}
			\end{bmatrix} \geq 0
		\end{equation}
		
		Используя S-процедуру, получаем:
		
		\begin{equation}
			\begin{bmatrix}
				(\bo{A} + \bo{B}\bo{K})\T\bo{S} + \bo{S} (\bo{A} + \bo{B}\bo{K}) + \gamma \bo{E}\T \bo{E} & \bo{S}\bo{F} \\
				\bo{F}\T\bo{S} & -\gamma\bo{I}
			\end{bmatrix} \prec 0
		\end{equation}		
		
		где $\gamma \geq 0$. Разделив полученное ЛМН на $\gamma$ и определив $\bo{S}_\gamma = \frac{1}{\gamma}\bo{S}$, получим:
		
		\begin{equation}
			\begin{bmatrix}
				(\bo{A} + \bo{B}\bo{K})\T\bo{S}_\gamma + \bo{S}_\gamma (\bo{A} + \bo{B}\bo{K}) + \bo{E}\T \bo{E} & \bo{S}_\gamma\bo{F} \\
				\bo{F}\T\bo{S}_\gamma & -\bo{I}
			\end{bmatrix} \prec 0
		\end{equation}
		
		
		
		
	\end{flushleft}
\end{frame}

\begin{frame}{Робастный синтез управления, 4}
	\begin{flushleft}
		
		Выражение:
		$$\begin{bmatrix}
			(\bo{A} + \bo{B}\bo{K})\T\bo{S}_\gamma + \bo{S}_\gamma (\bo{A} + \bo{B}\bo{K}) + \bo{E}\T \bo{E} & \bo{S}_\gamma\bo{F} \\
			\bo{F}\T\bo{S}_\gamma & -\bo{I}
		\end{bmatrix} \prec 0$$ 
		
		можно переписать как:
		
		\begin{equation*}
			\begin{bmatrix}
				(\bo{A} + \bo{B}\bo{K})\T\bo{S}_\gamma + \bo{S}_\gamma (\bo{A} + \bo{B}\bo{K}) & \bo{S}_\gamma\bo{F} \\
				\bo{F}\T\bo{S}_\gamma & -\bo{I}
			\end{bmatrix} 
			-
			\begin{bmatrix}
				\bo{E}\T \\
				0
			\end{bmatrix} 
			(-\bo{I})^{-1}
			\begin{bmatrix}
				\bo{E} & 0
			\end{bmatrix} 
			\prec 0
		\end{equation*}		
		
		затем, используя дополнение Шура, получаем:
		
		\begin{equation}
			\begin{bmatrix}
				(\bo{A} + \bo{B}\bo{K})\T\bo{S}_\gamma + \bo{S}_\gamma (\bo{A} + \bo{B}\bo{K}) & \bo{S}_\gamma\bo{F} & \bo{E}\T \\
				\bo{F}\T\bo{S}_\gamma & -\bo{I} & 0 \\
				\bo{E} & 0 & -\bo{I}
			\end{bmatrix} 
			\prec 0
		\end{equation}	
		
		
	\end{flushleft}
\end{frame}

\begin{frame}{Робастный синтез управления, 5}
	\begin{flushleft}
		
		Введём новую переменную $\bo{P} = \bo{S}^{-1}_\gamma$ и умножим обе части (конгруэнтное преобразование) ранее найденного неравенства на $\begin{bmatrix}
			\bo{P} & 0 & 0 \\
			0 & \bo{I} & 0 \\
			0 & 0 & \bo{I}
		\end{bmatrix}$, чтобы получить:
		
		\begin{equation}
			\begin{bmatrix}
				\bo{P}(\bo{A} + \bo{B}\bo{K})\T + (\bo{A} + \bo{B}\bo{K})\bo{P} & \bo{F} & \bo{P}\bo{E}\T \\
				\bo{F}\T & -\bo{I} & 0 \\
				\bo{E}\bo{P} & 0 & -\bo{I}
			\end{bmatrix} 
			\prec 0
		\end{equation}	
		
		Наконец, вводя последнюю новую переменную $\bo{L} = \bo{K}\bo{P}$, находим ЛМН, которое можно решить:
		
		\begin{equation}
			\begin{bmatrix}
				\bo{P}\bo{A}\T + \bo{A}\bo{P} + \bo{B}\bo{L} + \bo{L}\T \bo{B}\T & \bo{F} & \bo{P}\bo{E}\T \\
				\bo{F}\T & -\bo{I} & 0 \\
				\bo{E}\bo{P} & 0 & -\bo{I}
			\end{bmatrix} 
			\prec 0
		\end{equation}
		
		
	\end{flushleft}
\end{frame}

\begin{frame}{Квадратичная устойчивость, 1}
	\begin{flushleft}
		
		Рассмотрим следующую систему:
		
		\begin{equation}
			\dot{\bo{x}}  = \bo{A}\bo{x}
		\end{equation}
		%
		где $\bo{A} = \sum\limits_{i=1}^{n} \alpha_i \bo{A}_i$, $\alpha_i \geq 0$, $\sum\limits_{i=1}^{n} \alpha_i = 1$ с известными $\bo{A}_i$, но неизвестными коэффициентами $\alpha_i$. Будет ли система устойчива при всех возможных значениях $\alpha_i$? Обратите внимание, что в этом случае мы не можем использовать анализ собственных значений.
		
		\bigskip
		
		Геометрически это означает, что $\bo{A}$ находится в политопе с вершинами $\bo{A}_i$.
		
	\end{flushleft}
\end{frame}

\begin{frame}{Квадратичная устойчивость, 2}
	\begin{flushleft}
		
		\begin{theorem}[Квадратичная устойчивость]
			Если $\bo{A}_i^\top \bo{S} + \bo{S} \bo{A}_i \preceq 0$, то система $\dot{\bo{x}}  = \sum\limits_{i=1}^{n} \alpha_i \bo{A}_i \bo{x}$ устойчива, где $\alpha_i \geq 0$, $\sum\limits_{i=1}^{n} \alpha_i = 1$
		\end{theorem}
		
		\bigskip
		
		Доказательство: $\dot V = \left(\sum\limits_{i=1}^{n} \alpha_i \bo{A}_i \right)^\top \bo{S} + \bo{S} 
		\left( \sum\limits_{i=1}^{n} \alpha_i \bo{A}_i \right) \preceq 0$ можно переписать как: 
		$\dot V = \sum\limits_{i=1}^{n} \left( \alpha_i (\bo{A}_i^\top \bo{S} + \bo{S} \bo{A}_i) \right) $, и поскольку $\bo{A}_i^\top \bo{S} + \bo{S} \bo{A}_i \preceq 0$ и $\alpha_i \geq 0$, то $\dot V \leq 0$. \qed
		
	\end{flushleft}
\end{frame}

\begin{frame}{Квадратичная устойчивость - Синтез управления, 1}
	%	\framesubtitle{Часть 1}
	\begin{flushleft}
		
		Рассмотрим следующую систему:
		
		\begin{equation}
			\dot{\bo{x}}  = \bo{A}\bo{x} + \bo{B}\bo{u}
		\end{equation}
		%
		где $\bo{A} = \sum\limits_{i=1}^{n} \alpha_i \bo{A}_i$, $\alpha_i \geq 0$, $\sum\limits_{i=1}^{n} \alpha_i = 1$ с известными $\bo{A}_i$, но неизвестными коэффициентами $\alpha_i$. Как синтезировать закон управления $\bo{u} = \bo{K}\bo{x}$, делающий систему устойчивой при всех возможных значениях $\alpha_i$? 
		
		\bigskip
		
		Замкнутая форма системы имеет вид:
		
		\begin{equation}
			\dot{\bo{x}}  = \left(\sum\limits_{i=1}^{n} \alpha_i \bo{A}_i + \bo{B}\bo{K}\right)\bo{x}
		\end{equation}
		
		
	\end{flushleft}
\end{frame}

\begin{frame}{Квадратичная устойчивость - Синтез управления, 2}
	\begin{flushleft}
		
		Запишем уравнение Ляпунова для системы:
		
		\begin{equation}
			\left(
			\sum\limits_{i=1}^{n} \alpha_i (\bo{A}_i + \bo{B}\bo{K})
			\right)^\top \bo{S} 
			+ 
			\bo{S}
			\left(
			\sum\limits_{i=1}^{n} \alpha_i (\bo{A}_i + \bo{B}\bo{K})
			\right) 
			\prec 0
		\end{equation}
		
		Мы можем переписать его как:
		
		\begin{equation}
			\sum\limits_{i=1}^{n} \alpha_i 
			\left( 
			(\bo{A}_i + \bo{B}\bo{K})^\top \bo{S} +
			\bo{S} (\bo{A}_i + \bo{B}\bo{K})
			\right)
			\prec 0
		\end{equation}
		
		Следовательно, если $(\bo{A}_i + \bo{B}\bo{K})^\top \bo{S} +
		\bo{S} (\bo{A}_i + \bo{B}\bo{K}) \prec 0$, то исходная система устойчива.
		
	\end{flushleft}
\end{frame}

\begin{frame}{Квадратичная устойчивость - Синтез управления, 3}
	\begin{flushleft}
		
		Исходя из $(\bo{A}_i + \bo{B}\bo{K})^\top \bo{S} +
		\bo{S} (\bo{A}_i + \bo{B}\bo{K}) \prec 0$, можно перейти к синтезу управления. Вводя $\bo{P} = \bo{S}^{-1}$, используем конгруэнтное преобразование, умножая на $\bo{P}$ с обеих сторон:
		
		\begin{equation}
			\bo{P}(\bo{A}_i + \bo{B}\bo{K})^\top  +
			(\bo{A}_i + \bo{B}\bo{K})\bo{P} \prec 0
		\end{equation}
		
		Вводя новую переменную $\bo{L} = \bo{K} \bo{P}$, получаем задачу, линейную относительно переменных решения:
		
		\begin{equation}
			\bo{P}\bo{A}_i^\top + \bo{A}_i \bo{P} +
			\bo{L}^\top\bo{B}^\top + \bo{B}\bo{L} \prec 0
		\end{equation}		
		%
		где переменными решения являются $\bo{P}$ и $\bo{L}$. Матрица коэффициентов усиления управления находится как $\bo{K} = \bo{L} \bo{P}^{-1}$.
		
	\end{flushleft}
\end{frame}


\myqrframe



\begin{frame}{Appendix A}
	\framesubtitle{Congruence transformation and definiteness}
	\begin{flushleft}
		
		Consider matrices $\bo{P} \succeq 0$, and $\bo{V} \in \R^{n, n}$ is full rank. We can prove that:
		
		\begin{equation}
			\bo{P} \succ 0 \implies \bo{V}^\top\bo{P}\bo{V} \succ 0
		\end{equation}
		
		Proof: $\bo{x}^\top\bo{V}^\top\bo{P}\bo{V}\bo{x} = \bo{z}^\top\bo{P}\bo{z}$, where $\bo{z} = \bo{V}\bo{x}$. Since $\bo{P} \succ 0$, $\bo{z}^\top\bo{P}\bo{z} \geq 0$, hence $\bo{x}^\top\bo{V}^\top\bo{P}\bo{V}\bo{x} \geq 0$. 
		
		\begin{definition}
			Congruence transformation preserves semi-definiteness: $\text{det}(\bo{V}) \neq 0, \ \bo{P} \succ 0 \implies \bo{V}^\top\bo{P}\bo{V} \succ 0$
		\end{definition}
		
		
	\end{flushleft}
\end{frame}



\begin{frame}{Appendix B}
	\framesubtitle{S-procedure}
	\begin{flushleft}
		
There are two distinct results. The first, S-lemma, is non-conservative (necessary and sufficient condition):

\begin{theorem}[S-lemma]
	The following statements are equivalent. 
	
	\begin{enumerate}
		\item $\underset{\bo{z} \neq 0}{\text{sup}}\left\{ \bo{z}^\top \bo{M}\bo{z}  \ : \  \bo{z}^\top \bo{N}\bo{z} > 0 \right \} < 0$
		
		\item $\exists \gamma > 0$ such that $\bo{M} + \gamma \bo{N} \prec 0$
	\end{enumerate}
\end{theorem}

The second, conservative S-procedure, is a conservative (sufficient condition-only):

\begin{theorem}[Conservative S-procedure]
	If $\exists \gamma_i > 0$ such that $\bo{M} + \sum_{i} \gamma_i \bo{N}_i \preceq 0$, then $ \bo{z}^\top \bo{N}\bo{z} \geq 0 \implies \bo{z}^\top \bo{M}\bo{z}  \leq 0$
\end{theorem}
		
		
	\end{flushleft}
\end{frame}



\begin{frame}{Appendix B}
	\framesubtitle{S-procedure}
	\begin{flushleft}
		
		{\small 
		Earlier we proved the sufficiency of S-lemma. To prove the necessity, we consider the set of outputs of two quadratic forms: $f(\bo{z}) = \bo{z}^\top \bo{M}\bo{z}$ and $g(\bo{z}) = \bo{z}^\top \bo{N}\bo{z}$ defined as $\mathcal{C} = \{(g, f) \ : \ f = f(\bo{z}), \ g = g(\bo{z}), \ \forall \bo{z}\}$. $\mathcal{C}$ is a convex cone (joint range of two quadratic forms). 
		
		\bigskip
		
		Note that for the statement (1) of S-lemma implies that $\mathcal{C}$ intersects the positive cone $\mathcal{S} = \R^2_+$ only at the origin. Because both are convex cones, $\mathcal{C}$ and $\mathcal{S}$ can be separated by a hyperplane, for instance, a hyperplane with a normal vector $\bo{n} = (\mu, 1)$, where $\mu > 0$. For large $\mu$ the hyperplane aligns with the horizontal axis, for small - with the vertical one.
		
		\bigskip
		
		We can represent this separation as a condition on $\mathcal{C}$: $f + \mu g < 0$, which implies $\bo{z}^\top (\bo{M} + \gamma \bo{N})\bo{z} < 0$. This recovers the statement (2) of the lemma.
	}
		
	\end{flushleft}
\end{frame}



\begin{frame}{Schur Complement (Congruence Proof), 1}
	\begin{flushleft}
		
		Let 
		\[
		\bo{M} = 
		\begin{bmatrix}
			\bo{A} & \bo{B} \\
			\bo{B}^\top & \bo{C}
		\end{bmatrix},
		\quad \bo{C} = \bo{C}^\top \succ 0.
		\]
		
		We prove:
		\[
		\bo{M} \succeq 0 
		\quad \Longleftrightarrow \quad
		\bo{A} - \bo{B}\bo{C}^{-1}\bo{B}^\top \succeq 0.
		\]
		
		Define the congruence transformation:
		\[
		\bo{T} =
		\begin{bmatrix}
			\bo{I} & -\bo{B}\bo{C}^{-1} \\
			\bo{0} & \bo{I}
		\end{bmatrix}.
		\]
		
		Compute:
		\begin{align}
			\bo{T}^\top \bo{M} \bo{T}
			&=
			\begin{bmatrix}
				\bo{I} & \bo{0} \\
				-\bo{C}^{-1}\bo{B}^\top & \bo{I}
			\end{bmatrix}
			\begin{bmatrix}
				\bo{A} & \bo{B} \\
				\bo{B}^\top & \bo{C}
			\end{bmatrix}
			\begin{bmatrix}
				\bo{I} & -\bo{B}\bo{C}^{-1} \\
				\bo{0} & \bo{I}
			\end{bmatrix} \\
			&=
			\begin{bmatrix}
				\bo{A} - \bo{B}\bo{C}^{-1}\bo{B}^\top & \bo{0} \\
				\bo{0} & \bo{C}
			\end{bmatrix}.
		\end{align}
		
	\end{flushleft}
\end{frame}




\begin{frame}{Schur Complement (Congruence Proof), 2}
	\begin{flushleft}
		
		
		Congruence preserves semidefiniteness:
		\[
		\bo{M} \succeq 0 
		\;\Longleftrightarrow\;
		\bo{T}^\top \bo{M} \bo{T} \succeq 0
		\quad (\bo{T} \text{ invertible}).
		\]
		
		\bigskip
		
		Thus:
		\[
		\bo{M} \succeq 0
		\;\Longleftrightarrow\;
		\begin{bmatrix}
			\bo{A} - \bo{B}\bo{C}^{-1}\bo{B}^\top & 0 \\
			0 & \bo{C}
		\end{bmatrix} \succeq 0.
		\]
		
		\bigskip
		
		A block diagonal matrix is PSD iff each block is PSD:
		\[
		\Longleftrightarrow
		\begin{cases}
			\bo{A} - \bo{B}\bo{C}^{-1}\bo{B}^\top \succeq 0 \\
			\bo{C} \succeq 0
		\end{cases}
		\]
		
		\bigskip
		
		Since $\bo{C} \succ 0$, we conclude:
		\[
		\bo{M} \succeq 0 
		\;\Longleftrightarrow\;
		\bo{A} - \bo{B}\bo{C}^{-1}\bo{B}^\top \succeq 0. \ \ \ \qed
		\]
		
	\end{flushleft}
\end{frame}



\end{document}
